\chapter{Literature Review}
\label{ch:literature_review}

Introduction

The rapid expansion of biodiesel production has generated an oversupply of crude glycerol, a by-product that represents approximately 10\% (w/w) of the transesterification output (Garlapati et al., 2016; Attarbachi et al., 2023). This waste stream, loaded with impurities such as methanol, soaps, and salts, poses environmental disposal challenges but also offers a readily biodegradable carbon source for energy recovery (Anand \& Saxena, 2012; Gebreegziabher et al., 2025). Anaerobic digestion (AD) aligns with circular economy principles by converting organic residues into biogas while recycling nutrients, making crude glycerol a promising co-substrate (Gebreegziabher et al., 2025). In particular, AD in a CSTR, operated as a liquid-phase system, provides the continuous mixing and stable conditions needed for effective glycerol valorization (Li et al., 2011).

Glycerol as a Substrate for Anaerobic Digestion

Crude glycerol is highly biodegradable, and its co-digestion with nitrogen-rich substrates such as sewage sludge, animal manure, or food waste has repeatedly been shown to enhance methane yields. Gebreegziabher et al. (2025) reported that co-digestion with crude glycerol can boost methane production by 14–226\% compared to mono-digestion of the base substrate. Tomczak et al. (2025) reviewed one-stage AD experiments and concluded that a methane content above 60\% could be consistently obtained when the crude glycerol fraction did not exceed 8\% (v/v) of the feed. However, the same review emphasized that increasing the glycerol loading beyond a system-specific threshold may provoke inhibitory effects, likely due to the accumulation of volatile fatty acids (VFAs) and the impurities present in crude glycerol (Anand \& Saxena, 2012; Tomczak et al., 2025). The heterogeneous composition of crude glycerol—affected by oil feedstock, catalyst, and biodiesel processing—demands careful feedstock characterization and, if necessary, simple pretreatment to avoid process upsets (Attarbachi et al., 2023).

Continuous Stirred Tank Reactor (CSTR) Configuration

Liquid anaerobic digestion, typically conducted in CSTRs, operates at total solid concentrations of 0.5–15\% and is the dominant platform for treating slurries such as sewage sludge, manure, and food waste (Li et al., 2011). The CSTR’s continuous mixing provides uniform distribution of substrates, microorganisms, and heat, which is essential when feeding a rapidly fermentable co-substrate like glycerol. Temperature is a critical operational parameter; the activity of methanogenic consortia is strongly temperature-dependent, and both mesophilic and thermophilic regimes have been applied to AD (Choorit \& Wisarnwan, 2007). Tomczak et al. (2025) noted that the majority of glycerol co-digestion studies have been performed under mesophilic conditions, while thermophilic operation, which could offer higher conversion rates, still requires more systematic investigation. The CSTR platform is well suited to explore such temperature effects because its design allows easy control of heating and hydraulic retention time.

Role of Inoculum from Wastewater Treatment Plants

The selection of a robust microbial inoculum is decisive for rapid start-up and stable performance of a CSTR fed with glycerol. Anaerobic sludge from municipal wastewater treatment plants (WWTPs) is the most common source of inoculum because it harbours a diverse consortium of hydrolytic, acidogenic, and methanogenic microorganisms. Posmanik et al. (2020) demonstrated that granular sludge—a dense, biofilm-based biomass typically obtained from upflow anaerobic sludge blanket reactors treating wastewater—outperforms suspended biomass as an inoculum. In biochemical methane potential assays, granular sludge yielded higher ultimate methane production and exhibited markedly faster degradation kinetics for complex substrates. Although their work focused on batch tests, the findings strongly suggest that using granular sludge from a WWTP as seed for a CSTR could accelerate glycerol acclimation, improve process stability, and reduce the risk of acidification during load increases.

Process Stability and Monitoring

The easily degradable nature of glycerol makes the AD process susceptible to VFA accumulation and pH drop if the organic loading rate is not properly controlled. Therefore, close monitoring of process indicators is essential. Standard analytical protocols for water and wastewater (APHA, 2017) provide validated methods for measuring VFAs and bicarbonate alkalinity. Ribas et al. (2007) evaluated several methods for determining these parameters specifically in anaerobic reactors, highlighting the importance of accurate VFA and alkalinity measurements for timely detection of process imbalance. Maintaining a favourable VFA-to-alkalinity ratio is particularly critical when co-digesting glycerol, as excess VFAs can rapidly inhibit methanogenesis (Tomczak et al., 2025).

Conclusions

The available literature confirms that crude glycerol is a high-energy co-substrate that can significantly boost methane production in liquid AD systems. A CSTR provides the mixing and temperature control required to harness this potential. Using granular sludge from a wastewater treatment plant as the inoculum can enhance start-up speed and methanogenic activity, as demonstrated by comparative BMP studies. To unlock the full benefits while avoiding inhibition, the glycerol loading must be kept below documented thresholds, and the reactor state should be continuously tracked through standard measurements of VFAs and alkalinity. Further research is needed to optimise thermophilic CSTR operation and to establish long-term performance data for granular-sludge-inoculated reactors treating crude glycerol.
